% 2026-10-21 统一复核：几何、箭头、符号与正文逐项对应。
% 印宽17cm；最小字号8.4pt；保留独立注释行与框间空隙。
% 图契约：本题的源、边界和推导逐项对应；箭头符号按图注解释。
\documentclass[tikz,border=7pt]{standalone}
\usepackage{fontspec,xeCJK,amsmath,amssymb}
\IfFontExistsTF{Times New Roman}{\setmainfont{Times New Roman}}{\setmainfont{TeX Gyre Termes}}
\IfFontExistsTF{Heiti SC}{\setCJKmainfont[AutoFakeBold=1.5]{Heiti SC}}{\setCJKmainfont[AutoFakeBold=1.5]{FandolHei-Regular}}
\usetikzlibrary{arrows.meta,bending,calc,patterns,positioning}
\definecolor{ink}{HTML}{243745}
\definecolor{blue}{HTML}{275DAB}
\definecolor{green}{HTML}{177C63}
\definecolor{red}{HTML}{BD4935}
\definecolor{orange}{HTML}{AA6C16}
\definecolor{purple}{HTML}{715393}
\tikzset{every node/.style={font=\fontsize{9.3}{12.5}\selectfont,text=ink},
axis/.style={-{Stealth[length=4pt]},line width=.65pt,draw=ink!65},
vec/.style={-{Stealth[length=5pt]},line width=1pt,draw=blue},
short/.style={-{Stealth[length=3pt]},line width=.8pt},
guide/.style={dashed,draw=ink!40,line width=.5pt},
card/.style={rounded corners=2pt,draw=blue!65,fill=blue!4,align=center,text width=6.5cm,inner xsep=9pt,inner ysep=7pt},
title/.style={font=\fontsize{10.5}{14}\selectfont\bfseries,text width=15cm,align=center},
label/.style={font=\fontsize{8.4}{11}\selectfont,align=center},
note/.style={label,text width=6.2cm}}
\begin{document}\begin{tikzpicture}[x=1cm,y=1cm]
\path[use as bounding box] (-8,-3.5) rectangle(8,3.5);

\node[title] at(0,3.15){有源区与无源区：先划区域，再写方程};
\begin{scope}[shift={(-4.25,.15)}]
\filldraw[draw=blue!70,fill=blue!4] (0,0)circle(2);
\filldraw[draw=red!75,fill=red!8] (0,0)circle(.95);
\node[align=center] at(0,0){体电荷区\\$\rho\ne0$};
\node[align=center] at(0,1.45){无源区 $\rho=0$};
\node[label] at(0,-2.45){真空；界面上无面电荷};
\end{scope}
\node[card] at(3.6,1.65){有源区：$\displaystyle\nabla^2\varphi=-\rho/\varepsilon_0$};
\node[card,draw=green!70,fill=green!4] at(3.6,.05){无源区：$\nabla^2\varphi=0$};
\node[card,draw=orange!70,fill=orange!4] at(3.6,-1.55){在无面电荷的真空界面\\连接 $\varphi$ 与法向导数 $\partial_n\varphi$};
\node[label,text width=15.1cm] at(0,-3.05){“无源”指该区域内总电荷密度为零；不能跨过点电荷或面电荷，直接写同一个拉普拉斯方程。};
\end{tikzpicture}\end{document}
